%%%%%%%%%%%%%%%%%%%%%%%%%%%%%%%%%%%%%%%%%%%%%%%%%%%%%%%%%%%%%%%%%%%%%%%%%%%%%%%
% APPENDIX G: GLOSSARY AND NOTATION
% Complementarity and Context Framework
% Version: v53.0
% Last updated: 2026-01-06 (v53: added 5C CORE, BBB concepts, axiom index)
%%%%%%%%%%%%%%%%%%%%%%%%%%%%%%%%%%%%%%%%%%%%%%%%%%%%%%%%%%%%%%%%%%%%%%%%%%%%%%%

\section{Glossary and Notation}
\label{app:glossary}

%==============================================================================
% HEADER BLOCK (Template Compliance)
%==============================================================================

\begin{tcolorbox}[colback=blue!5!white, colframe=blue!75!black,
    title=\textbf{Appendix G: REF-GLOSSARY}]

\begin{tabular}{@{}ll@{}}
\textbf{Appendix:} & G (Central Glossary and Notation) \\
\textbf{Category:} & REF (Reference Materials) \\
\textbf{Scope:} & Comprehensive symbol, term, and axiom definitions for EBF \\
\textbf{Version:} & 53.0 \\
\textbf{Last Updated:} & January 2026 \\
\textbf{Dependencies:} & All EBF appendices (this is the central reference) \\
\textbf{Outputs:} & Standardized notation across all chapters and appendices \\
                  & Complete axiom index \\
                  & Cross-reference navigation map \\
\end{tabular}

\end{tcolorbox}

%==============================================================================
% ABSTRACT (Template Compliance)
%==============================================================================

\begin{tcolorbox}[colback=white, colframe=black,
    title=\textbf{Abstract}]

This appendix serves as the central reference for all symbols, terms, and concepts
in the EBF framework. It provides: (1) comprehensive mathematical notation for
complementarity structures ($C$, $C^*$), context dimensions ($\Psi$), and welfare
measures ($K$, $Q$, $W$); (2) definitions of the FEPSDE welfare decomposition and
INU/KNU/IDN utility categories; (3) complete axiom indices for CORE appendices
(V, BBB, AU, AV); and (4) cross-reference navigation to detailed treatments.
This glossary ensures terminological consistency across all 50+ appendices and
15 chapters of the EBF documentation.

\textit{(85 words)}

\end{tcolorbox}

%==============================================================================
% QUICK REFERENCE BOX
%==============================================================================

\begin{tcolorbox}[colback=yellow!10!white, colframe=yellow!75!black,
    title=\textbf{Quick Reference: Core Symbols}]

\textbf{Complementarity:}
$C$ = actual structure \quad $C^*$ = reference structure \quad $\Gamma$ = complementarity matrix

\textbf{Measures:}
$K = 1 - \|C - C^*\|/\|C^*\|$ (coherence) \quad $Q$ = quality \quad $W = Q \cdot g(K)$ (welfare)

\textbf{Context:}
$\Psi = (\Psi_I, \Psi_S, \Psi_C, \Psi_K, \Psi_E, \Psi_T, \Psi_M, \Psi_F)$ (8 dimensions)

\textbf{Utility:}
INU (Individual) + KNU (Collective) + IDN (Identity) $\times$ FEPSDE (6 dimensions)

\textbf{Awareness \& Willingness:}
$U_{eff} = A(\cdot) \times U_{pot}$ \quad Action $\Leftrightarrow WAX \geq \theta$

\end{tcolorbox}

%==============================================================================
% CHAPTER LINKAGE
%==============================================================================

\begin{tcolorbox}[colback=purple!5!white, colframe=purple!75!black,
    title=\textbf{Chapter Linkages}]

\textbf{Primary Role:} Central reference for all chapters and appendices

\textbf{Key Integrations:}
\begin{itemize}[nosep]
    \item All chapters reference this glossary for notation consistency
    \item CORE appendices (AAA, B, C, V, BBB, AU, AV) define symbols used here
    \item FORMAL appendices (A, D) provide mathematical derivations
    \item All other appendices follow notation defined here
\end{itemize}

\textbf{Usage Guide:}
\begin{itemize}[nosep]
    \item For symbol lookup: See Section \ref{sec:glossary-symbols} (Core Mathematical Symbols)
    \item For $\Psi$ dimensions: See detailed definitions below
    \item For axiom references: See Section \ref{sec:glossary-axioms}
\end{itemize}

\end{tcolorbox}

%==============================================================================
% CROSS-REFERENCE MAP
%==============================================================================

\begin{tcolorbox}[colback=white, colframe=black,
    title=\textbf{Cross-Reference Map}]

\begin{center}
\small
\begin{tabular}{@{}c|cccccc@{}}
\toprule
& \textbf{AAA} & \textbf{B} & \textbf{C} & \textbf{V} & \textbf{BBB} & \textbf{AU/AV} \\
\midrule
\textbf{G $\leftarrow$} & Levels $L$ & $\Gamma$-matrix & FEPSDE & $\Psi$ dims & Estimation & $A(\cdot)$, WAX \\
\textbf{G $\rightarrow$} & Notation & Notation & Notation & Notation & Notation & Notation \\
\midrule
\textbf{Link} & \cellcolor{red!20}Strong & \cellcolor{red!20}Strong & \cellcolor{red!20}Strong & \cellcolor{red!20}Strong & \cellcolor{red!20}Strong & \cellcolor{red!20}Strong \\
\bottomrule
\end{tabular}
\end{center}

\textbf{Legend:}
\colorbox{red!20}{Strong} = Bidirectional dependency (all appendices reference G)

\end{tcolorbox}

This appendix provides a comprehensive reference for all symbols, terms, and concepts used throughout the framework. Updated 2026-01-06 (v52.1: corrected FEPSDE 144-component structure to match Appendix C).

%==============================================================================
% \section{Theory} marker for template compliance
\subsection{Theoretical Foundation: Notation Principles}
\label{sec:glossary-theory}
%==============================================================================

The EBF notation system follows three principles:

\begin{enumerate}
    \item \textbf{Consistency:} Same symbol means same concept across all documents
    \item \textbf{Hierarchy:} Greek letters for parameters, Latin for variables, calligraphic for sets
    \item \textbf{Context-dependence:} Parenthetical $(\Psi)$ indicates context modulation
\end{enumerate}

%==============================================================================
\subsection{Core Mathematical Symbols}
\label{sec:glossary-symbols}
%==============================================================================

\begin{center}
\renewcommand{\arraystretch}{1.3}
\begin{longtable}{lp{4cm}p{7cm}}
\hline
\textbf{Symbol} & \textbf{Name} & \textbf{Definition} \\
\hline
\endfirsthead
\hline
\textbf{Symbol} & \textbf{Name} & \textbf{Definition} \\
\hline
\endhead
\hline
\endfoot

\multicolumn{3}{l}{\textbf{Complementarity Structure}} \\
\hline
$C$ & Complementarity matrix & Cross-partial derivatives: $C_{ij} = \partial^2 U / \partial x_i \partial x_j$ \\
$C^*$ & Reference structure & Self-consistent fixed point satisfying $C^* = \Phi(C^*, \Psi)$ \\
$C^*(\Psi)$ & Context-dependent reference & Reference structure as function of context \\
$C_{ij}$ & Complementarity coefficient & Strength of interaction between components $i$ and $j$ \\
$\hat{C}$ & Estimated complementarity & Empirically calibrated complementarity matrix \\
$\Delta$ & Deviation matrix & $\Delta = C - C^*$; diagnostic signal for learning \\
$\|\Delta\|$ & Deviation magnitude & Frobenius norm of deviation matrix \\
\hline

\multicolumn{3}{l}{\textbf{Quality and Coherence Measures}} \\
\hline
$K$ & Coherence index & $K = 1 - \|C - C^*\|/\|C^*\|$; alignment with reference \\
$K^*$ & Optimal coherence & Maximum achievable coherence given constraints \\
$Q$ & Quality index & Welfare potential at equilibrium \\
$W$ & Realized welfare & $W = Q \cdot g(K)$; actual welfare achieved \\
$g(K)$ & Coherence function & Maps coherence to welfare realization \\
$\kappa$ & Coherence threshold & Critical $K$ value for regime transition \\
\hline

\multicolumn{3}{l}{\textbf{Context Variables}} \\
\hline
$\Psi$ & Context vector & Eight-dimensional vector $(\Psi_I, \Psi_S, \Psi_C, \Psi_K, \Psi_E, \Psi_T, \Psi_M, \Psi_F)$ \\
$\Psi_I$ & Institutional context & Rules, enforcement, property rights strength \\
$\Psi_S$ & Social context & Trust, norms, social capital \\
$\Psi_C$ & Cognitive context & Information processing capacity, numeracy \\
$\Psi_K$ & Informational context & Transparency, signal quality, asymmetries \\
$\Psi_E$ & Economic context & Development level, market thickness \\
$\Psi_T$ & Temporal context & Time horizons, stability, uncertainty \\
$\Psi_M$ & Market scope & Geographic reach, trade openness \\
$\Psi_F$ & Factor flexibility & Adjustment capacity, mobility \\
$\bar{\Psi}$ & Reference context & Baseline context for comparison \\
$\Delta\Psi$ & Context shift & Change in context: $\Delta\Psi = \Psi' - \Psi$ \\
\hline

\multicolumn{3}{l}{\textbf{Behavioral Parameters}} \\
\hline
$\lambda$ & Loss aversion & Multiplier for losses relative to gains \\
$\lambda_d$ & Domain-specific loss aversion & Loss aversion in FEPSDE dimension $d$ \\
$\lambda(\Psi)$ & Context-dependent loss aversion & $\lambda = f(\Psi_C, \Psi_E, \Psi_T, ...)$ \\
$\delta$ & Discount factor & Time preference parameter \\
$\delta_d$ & Domain-specific discount & Discount factor in dimension $d$ \\
$\beta$ & Present bias & Quasi-hyperbolic discounting parameter \\
$\alpha$ & Reference point weight & Importance of reference in valuation \\
$\gamma$ & Misfit penalty & Cost of deviation from reference \\
$\gamma(\Psi)$ & Context-dependent misfit & Penalty varies with institutional context \\
$\eta$ & Curvature parameter & Risk attitude in prospect theory \\
$\theta$ & Probability weighting & Distortion of objective probabilities \\
\hline

\multicolumn{3}{l}{\textbf{Social Preferences}} \\
\hline
$\alpha_F$ & Fairness concern & Weight on fairness considerations \\
$\beta_R$ & Reciprocity parameter & Strength of reciprocal responses \\
$\rho$ & Inequality aversion & Sensitivity to distributional inequality \\
$\sigma_S$ & Social signaling & Weight on reputation/status \\
$\tau$ & Trust parameter & Baseline trust in social interactions \\
\hline

\multicolumn{3}{l}{\textbf{Strategic Variables}} \\
\hline
$s$ & Strategy axis & Goal orientation $\in [0,1]$; efficiency vs. legitimacy \\
$\sigma$ & Structure axis & Organizational mode $\in [0,1]$; hierarchy vs. network \\
$(s^*, \sigma^*)$ & Optimal configuration & Strategy-structure pair maximizing $W$ \\
$\Pi$ & Payoff function & Returns to strategic choices \\
\hline

\multicolumn{3}{l}{\textbf{Welfare Components}} \\
\hline
$U$ & Utility function & Individual welfare representation \\
$U^{INU}$ & Individual utility & Personal welfare component \\
$U^{KNU}$ & Collective utility & Social/group welfare component \\
$U^{IDN}$ & Identity utility & Self-concept and meaning component \\
$V$ & Value function & Reference-dependent valuation \\
$W^{FEPSDE}$ & FEPSDE welfare & Six-dimensional welfare measure \\
\hline

\multicolumn{3}{l}{\textbf{Dynamics and Transitions}} \\
\hline
$\dot{C}$ & Complementarity dynamics & Time derivative of $C$ \\
$T(\Psi, \Psi')$ & Transition cost & Cost of moving from $\Psi$ to $\Psi'$ \\
$v$ & Adjustment velocity & Speed of convergence to $C^*$ \\
$\tau_{adj}$ & Adjustment time & Time to reach new equilibrium \\
$p_{trans}$ & Transition probability & Likelihood of regime change \\
\hline

\multicolumn{3}{l}{\textbf{Estimation and Calibration}} \\
\hline
$\hat{\theta}$ & Parameter estimate & Calibrated parameter value \\
$SE(\hat{\theta})$ & Standard error & Uncertainty in estimate \\
$N$ & Sample size & Number of observations/experiments \\
$R^2$ & Coefficient of determination & Variance explained \\
$\chi^2$ & Chi-squared statistic & Goodness of fit measure \\
\hline

\multicolumn{3}{l}{\textbf{Segmentation and Intervention (Axioms 8-13)}} \\
\hline
$C_k$ & Behavioral cluster $k$ & Natural segment of population with homogeneous preferences \\
$k$ & Cluster count & Number of natural segments (context-dependent) \\
$k^*$ & Optimal cluster count & $\arg\min_k \text{BIC}$ or optimal silhouette score \\
$\mathcal{L}_k$ & Leverage profile set & Collection of leverage dimensions for cluster $k$ \\
$\ell_k^d$ & Leverage in dimension $d$ & $\text{Var}(\omega_i^d | i \in C_k) / \sum_{d'} \text{Var}(\omega_i^{d'} | i \in C_k)$ \\
$\mathcal{D}_I$ & Intervention design set & Collection of intervention intensities across dimensions \\
$\delta_I^d$ & Intervention intensity & Strength of intervention targeting dimension $d \in [0,1]$ \\
$I^*_k$ & Optimized intervention & Segment-specific intervention maximizing leverage match \\
$\text{Leverage}(I, C_k)$ & Matching function & $\sum_d (\ell_k^d \cdot \delta_I^d) + \sum_b (\gamma_k^b \cdot \beta_I^b)$ \\
$\text{AvgLeverage}$ & Average leverage match & Mean matching quality across all clusters \\
$E[\Delta O_{\text{seg}}]$ & Expected segment effect & Predicted outcome improvement from segmentation \\
$E[\Delta O_{\text{non-seg}}]$ & Expected population effect & Baseline outcome improvement without segmentation \\
$\beta^{\text{match}}$ & Matching coefficient & Effect size scaling with leverage match \\
$\omega_i^d$ & Preference in dimension $d$ & Individual $i$'s utility weight on dimension $d \in \{F,E,P,S,D,E\}$ \\
\hline

\end{longtable}
\end{center}

%==============================================================================
\subsection{The Eight $\Psi$ Dimensions: Detailed Definitions}
%==============================================================================

\begin{center}
\renewcommand{\arraystretch}{1.3}
\begin{longtable}{lp{3cm}p{5cm}p{3cm}}
\hline
\textbf{Dim.} & \textbf{Name} & \textbf{Components} & \textbf{Measurement} \\
\hline
\endfirsthead
\hline
\textbf{Dim.} & \textbf{Name} & \textbf{Components} & \textbf{Measurement} \\
\hline
\endhead
\hline
\endfoot

$\Psi_I$ & Institutional & Rule of law, contract enforcement, property rights, regulatory quality, corruption control & WGI indicators, PRS scores \\

$\Psi_S$ & Social & Generalized trust, social capital, norm compliance, community ties, reciprocity norms & WVS trust items, GSS social capital \\

$\Psi_C$ & Cognitive & Numeracy, financial literacy, cognitive reflection, working memory, decision quality & CRT scores, numeracy tests, PISA \\

$\Psi_K$ & Informational & Information availability, signal quality, transparency, media freedom, asymmetry reduction & Press freedom index, disclosure scores \\

$\Psi_E$ & Economic & GDP per capita, market development, financial depth, infrastructure quality, resource access & GDP/capita, market cap/GDP \\

$\Psi_T$ & Temporal & Planning horizon, patience, stability expectations, uncertainty tolerance, future orientation & Discount rate elicitation, stability indices \\

$\Psi_M$ & Market Scope & Trade openness, geographic market access, export diversity, integration depth & Trade/GDP, Herfindahl indices \\

$\Psi_F$ & Factor Flexibility & Labor mobility, capital reallocation, skill transferability, adjustment speed & Job flow rates, retraining data \\
\hline

\end{longtable}
\end{center}

%==============================================================================
\subsection{FEPSDE Welfare Dimensions}
%==============================================================================

The FEPSDE framework decomposes welfare into six fundamental dimensions:

\begin{center}
\renewcommand{\arraystretch}{1.3}
\begin{tabular}{llp{8cm}}
\hline
\textbf{Dim.} & \textbf{Name} & \textbf{Definition and Components} \\
\hline
F & Financial & Monetary wealth, income, consumption, economic security, asset ownership \\
E & Emotional & Subjective well-being, affect balance, life satisfaction, mental health \\
P & Physical & Health status, longevity, bodily functioning, energy, physical safety \\
S & Social & Relationships, belonging, status, recognition, community integration \\
D & Digital & Data rights, algorithmic treatment, platform access, digital autonomy \\
Eco & Ecological & Environmental quality, sustainability, climate security, natural capital \\
\hline
\end{tabular}
\end{center}

The complete FEPSDE matrix has $3 \times 2 \times 6 \times 4 = 144$ components:
\begin{itemize}
    \item 3 utility categories: INU (Individual), KNU (Collective), IDN (Identity)
    \item 2 valence types: Gain, Pain
    \item 6 FEPSDE dimensions: F, E, P, S, D, Eco
    \item 4 time horizons: Instant, Short, Medium, Long
\end{itemize}

%==============================================================================
\subsection{Utility Categories (EBF Architecture)}
%==============================================================================

\begin{center}
\renewcommand{\arraystretch}{1.3}
\begin{tabular}{llp{7cm}}
\hline
\textbf{Code} & \textbf{Name} & \textbf{Definition} \\
\hline
INU & Individual Utility & Personal welfare from own consumption, health, experiences \\
KNU & Collective Utility & Welfare from group outcomes, public goods, social welfare \\
IDN & Identity Utility & Welfare from self-concept, meaning, values alignment \\
\hline
AWX & Awareness Module & Information processing, attention, belief formation \\
WIL & Willingness Module & Motivation, effort, behavioral activation \\
\hline
\end{tabular}
\end{center}

%==============================================================================
\subsection{Key Concepts and Terms: Intuitive and Formal Definitions}
\label{sec:glossary-concepts}
%==============================================================================

Each concept is defined both \textbf{intuitively} (accessible explanation with examples) and \textbf{formally} (mathematical specification). This dual approach ensures accessibility for practitioners while maintaining rigor for researchers.

\begin{center}
\small
\renewcommand{\arraystretch}{1.5}
\begin{longtable}{@{}p{2.5cm}p{5.2cm}p{5.2cm}@{}}
\toprule
\textbf{Begriff} & \textbf{Intuitive Definition} & \textbf{Formale Definition} \\
\midrule
\endfirsthead
\toprule
\textbf{Begriff} & \textbf{Intuitive Definition} & \textbf{Formale Definition} \\
\midrule
\endhead
\midrule
\multicolumn{3}{r}{\textit{Fortsetzung nächste Seite...}} \\
\bottomrule
\endfoot
\bottomrule
\endlastfoot

\multicolumn{3}{l}{\cellcolor{blue!10}\textbf{Kernkonzepte (Core Concepts)}} \\
\midrule

Complementarity &
Wie stark beeinflusst das Verhalten einer Person den Nutzen einer anderen? Wenn Ihr Kollege kooperiert, wird Ihre eigene Kooperation wertvoller---das ist positive Komplementarität. &
$C_{ij} = \frac{\partial^2 U_i}{\partial x_i \partial x_j}$; \newline $C_{ij} > 0$: Komplemente \newline $C_{ij} < 0$: Substitute \newline $C_{ij} = 0$: unabhängig \\

Context &
Alles, was die ``Spielregeln'' bestimmt, aber nicht von Einzelnen kontrolliert wird: Gesetze, Kultur, Technologie, Vertrauen. Der Kontext bestimmt, welches Verhalten optimal ist. &
$\Psi = (\Psi_I, \Psi_S, \Psi_C, \Psi_K, \Psi_E, \Psi_T, \Psi_M, \Psi_F)$ \newline 8 Dimensionen, je $\in [0,1]$ \\

Reference Structure &
Das ``ideale'' Verhaltensmuster für einen gegebenen Kontext. Wie ein Navi: zeigt den optimalen Weg, abhängig davon, wo Sie sind. &
$C^*(\Psi)$: Fixpunkt mit \newline $C^* = \Phi(C^*, \Psi)$ \\

Coherence &
Wie gut passt tatsächliches Verhalten zur Referenz? Hohe Kohärenz = Menschen verhalten sich kontextangemessen. &
$K = 1 - \frac{\|C - C^*\|}{\|C^*\|}$ \newline $K \in [0,1]$; $K=1$: perfekt \\

\midrule
\multicolumn{3}{l}{\cellcolor{green!10}\textbf{Welfare und Utility}} \\
\midrule

Utility &
Was Menschen wertvoll finden---nicht nur Geld, sondern auch Gesundheit, Beziehungen, Identität. Drei Arten: Eigen-, Gemeinschafts-, Identitätsnutzen. &
$U = U^{INU} + U^{KNU} + U^{IDN}$ \newline über FEPSDE-Dimensionen \\

Quality Index &
Das Wohlfahrtspotential bei optimalem Verhalten. Wie der theoretische Verbrauch eines Autos unter Idealbedingungen. &
$Q = U(C^*, \Psi)$ \\

Welfare &
Der tatsächlich realisierte Nutzen, abhängig von der Kohärenz. Realität ist meist niedriger als Potential. &
$W = Q \cdot g(K)$; \newline $W \leq Q$ \\

FEPSDE &
Die sechs Dimensionen menschlichen Wohlbefindens: Financial, Emotional, Physical, Social, Digital, Ecological. &
$6 \times 3 \times 2 \times 4 = 144$ Komponenten \\

\midrule
\multicolumn{3}{l}{\cellcolor{yellow!10}\textbf{Behavioral Parameters}} \\
\midrule

Loss Aversion &
Verluste schmerzen mehr als gleich große Gewinne erfreuen. 100€ verlieren fühlt sich schlimmer an als 100€ gewinnen gut. &
$\lambda > 1$; typisch $\lambda \approx 2$; \newline $\lambda = f(\Psi_C, \Psi_E, ...)$ \\

Present Bias &
Die Tendenz, sofortige Belohnungen überzubewerten. ``Lieber heute 10€ als morgen 11€.'' &
$\beta < 1$ in quasi-hyperbolischem Diskontieren \\

Social Preferences &
Menschen kümmern sich um andere---Fairness, Reziprozität. Wir lehnen unfaire Angebote ab, auch wenn es uns kostet. &
$U_i = u_i + \alpha_F \cdot \text{Fair} + \beta_R \cdot \text{Rez}$ \\

Reference Point &
Der Maßstab für Bewertungen. Ob eine Gehaltserhöhung ``gut'' ist, hängt von Erwartungen ab. &
$V(x) = v(x - r)$; \newline $r$ = Referenzpunkt \\

Endowment Effect &
Was wir besitzen, bewerten wir höher als Äquivalentes, das wir nicht besitzen. ``Meine Tasse ist 10€ wert, deine nur 5€.'' &
Spezialfall von $\lambda > 1$; \newline Besitz verschiebt $r$ \\

\midrule
\multicolumn{3}{l}{\cellcolor{red!10}\textbf{9C CORE Framework}} \\
\midrule

Awareness &
Was bemerken Menschen überhaupt? Nicht alle relevanten Informationen werden wahrgenommen. Awareness filtert Entscheidungsinput. &
$U_{eff} = A(\cdot) \times U_{pot}$; \newline $A: [0,1]$ \\

Willingness &
Selbst wenn Menschen wissen, was optimal wäre---sind sie bereit zu handeln? Motivation und Handlungsbereitschaft. &
Aktion $\Leftrightarrow WAX \geq \theta$; \newline WAX = Willingness-Index \\

Stage &
Wo steht jemand auf der Verhaltensänderungsreise? Von ``Ich weiß nicht, dass ich ein Problem habe'' bis ``dauerhaft geändert''. &
$S(t) \in \{$Pre, Cont, Prep, Act, Maint$\}$ \\

\midrule
\multicolumn{3}{l}{\cellcolor{purple!10}\textbf{Methodik}} \\
\midrule

Calibration &
Parameter aus echten Experimenten schätzen. ``Was ist $\lambda$ wirklich?'' basierend auf beobachtetem Verhalten. &
$\hat{\theta} = \arg\max \mathcal{L}(\text{data}|\theta)$ \\

LLM Monte Carlo &
KI-gestützte Methode zur Parameterschätzung. LLMs generieren Szenarien; Statistik extrahiert Parameter. &
Appendix AN: $\hat{C}_{ij} = \mathbb{E}_{LLM}[\cdot]$ \\

Metatheory &
Framework, das Theorien organisiert, nicht ersetzt. Sagt, \emph{wann} welche Theorie gilt. &
$\mathcal{T}(\Psi) \mapsto \{$Arrow-Debreu, Behavioral, ...$\}$ \\

Digital Twin &
KI-Simulation eines Menschen. Aktuelle Forschung: repliziert nur 50\% der Experimente korrekt. &
Peng et al.\ (2025): 50\% Replikation \\

\midrule
\multicolumn{3}{l}{\cellcolor{gray!10}\textbf{Dynamics und Transitions}} \\
\midrule

Coherence Trap &
Hohe interne Konsistenz verdeckt Veränderungsbedarf. ``Wir machen alles richtig''---aber der Kontext hat sich geändert. Wie Kodak. &
$K$ hoch aber $\Psi$ verschiebt; \newline $\partial C^*/\partial \Psi$ ignoriert \\

Regime Transition &
Diskrete Sprünge wenn Kontext kritische Schwellen überquert. Fundamentaler Strukturwandel, nicht graduelle Anpassung. &
$\Psi > \bar{\Psi}_{crit}$ $\Rightarrow$ $C^*_{new}$ \\

Transformation &
Der Prozess des Strukturwandels. Kostspielig, zeitaufwändig, aber manchmal notwendig. &
$T = \int c(C, \dot{C}) dt$ \\

Treatment Effect Heterogeneity &
Interventionen wirken unterschiedlich je nach Kontext. Was in Schweden funktioniert, kann in Nigeria scheitern. &
$\frac{\partial \text{ATE}}{\partial \Psi} \neq 0$ \\

Functional Knowledge &
Nicht nur ``Verlustaversion existiert'', sondern ``$\lambda = 2.1$ wenn $\Psi_C = 0.7$''. Quantitative Beziehungen. &
$\lambda(\Psi) = f(\Psi_C, \Psi_E, ...)$ \\

Zürich Tradition &
Die experimentelle Wirtschaftsforschung von Ernst Fehr und Kollegen: Fairness, Vertrauen, Reziprozität, soziale Präferenzen. &
Appendix K: 150+ Referenzen; \newline Kapitel 5 \\

\midrule
\multicolumn{3}{l}{\cellcolor{orange!10}\textbf{Segmentation and Intervention Matching}} \\
\midrule

Behavioral Segment &
Natürliche Gruppen in der Bevölkerung mit ähnlichen Präferenzen und Verhaltensmustern. Nicht künstlich definiert, sondern datengesteuert aus Heterogenität entdeckt. &
$C_k$: Cluster mit homogenen $(\omega_i, A_i, W_i)$ \newline $k^* = \arg\min_k \text{BIC}$ \\

Leverage Profile &
Welche Motivationsdimensionen (FEPSDE) treiben das Verhalten in einem Segment? Identifiziert die ``Hebelpunkte'' für gezielte Interventionen. &
$\ell_k^d = \frac{\text{Var}(\omega_i^d | i \in C_k)}{\sum_{d'} \text{Var}(\omega_i^{d'} | i \in C_k)}$ \\

Intervention Design &
Wie intensiv sollte eine Intervention jede Motivationsdimension addressieren? Bindet sich auf 1.0 normalisierte Intensität je Dimension. &
$\delta_I^d \in [0,1]$: Intensität von Intervention $I$ auf Dimension $d$ \\

Intervention Matching &
Gute Segmentationsmessungen produzieren nur Wert wenn Interventionen auf Leverage abgestimmt sind. Hohe Leverage + schlecht designte Interventionen = keine Effekte. &
$\text{Leverage}(I, C_k) = \sum_d (\ell_k^d \cdot \delta_I^d)$ \\

Theory Externalization &
Wenn Praktiker ein Experiment designen, müssen sie ihre impliziten Intuition in explizite mathematische Strukturen konvertieren. Kraft dieses Prozesses macht Wissen transferierbar. &
Spezifiziere $\mathcal{L}_k$ und $\mathcal{D}_I$ vor Datenerfassung (Axiom 10) \\

Pre-Registration &
Bevor Ergebnisse gesehen werden, festlegen: Primäre Vorhersagen, Effektgrößen, Diagnostisches Entscheidungsbaum. Verhindert Post-Hoc Rationalisierung. &
Pre-register: $E[\Delta O_{\text{seg}}] = \beta^{\text{match}} \cdot \text{AvgLeverage}$ (Axiom 11) \\

Diagnostic Refinement &
Ergebnisse sind nicht nur ``Ja/Nein''. Sie diagnostizieren wo genau Theorie fehlgeschlagen hat: schlechte Segmentierung oder schlechtes Intervention-Matching? &
Berechne $\text{AvgLeverage}_{\text{observed}}$ um Problem zu lokalisieren (Axiom 12) \\

Cumulative Science &
Der Zyklus: externalisieren → testen → diagnostizieren → verfeinern → transferieren. Erzeugt Wissen das sich ansammelt, nicht Handwerk das verloren geht. &
Registry FFF speichert: $(\mathcal{L}_k, \mathcal{D}_I, \beta^{\text{match}}, \text{Erfolg/Fehler})$ (Axiom 13) \\

\end{longtable}
\end{center}

%==============================================================================
\subsection{Acronyms and Abbreviations}
%==============================================================================

\begin{center}
\renewcommand{\arraystretch}{1.2}
\begin{longtable}{lp{10cm}}
\hline
\textbf{Acronym} & \textbf{Meaning} \\
\hline
\endfirsthead
\hline
\textbf{Acronym} & \textbf{Meaning} \\
\hline
\endhead
\hline
\endfoot

AWX & Awareness-Complementarity Module \\
ATE & Average Treatment Effect \\
BIC & Bayesian Information Criterion \\
BCM & Behavioral Change Model \\
EBF & Evidence-Based Framework for Economic and Social Behavior \\
BEATRIX & Behavioral Economics AI-Integrated Research and Intervention eXpert \\
CRT & Cognitive Reflection Test \\
DG & Dictator Game \\
FEPSDE & Financial, Emotional, Physical, Social, Digital, Ecological \\
FOC & First-Order Condition \\
GSS & General Social Survey \\
IDN & Identity Utility \\
INU & Individual Utility \\
KNU & Collective Utility (German: Kollektiver Nutzen) \\
LLM & Large Language Model \\
MPC & Marginal Propensity to Consume \\
OLS & Ordinary Least Squares \\
PD & Prisoner's Dilemma \\
PG & Public Goods Game \\
PISA & Programme for International Student Assessment \\
PRS & Political Risk Services \\
PT & Prospect Theory \\
RCT & Randomized Controlled Trial \\
SE & Standard Error \\
SWB & Subjective Well-Being \\
TEH & Treatment Effect Heterogeneity \\
FFF & Framework Falsification \& Feedback Registry (Appendix FFF: METHOD-REGISTRY) \\
GGG & Goal Guidance \& Generalization Tables (Appendix GGG: METHOD-CONFIG) \\
EEE & Externalization Experimentation \& Enhancement Workflow (Appendix EEE: METHOD-DESIGN) \\
UG & Ultimatum Game \\
WGI & World Governance Indicators \\
WIL & Willingness Module \\
WVS & World Values Survey \\
\hline

\end{longtable}
\end{center}

%==============================================================================
\subsection{Mathematical Notation Conventions}
%==============================================================================

\begin{center}
\renewcommand{\arraystretch}{1.3}
\begin{tabular}{lp{10cm}}
\hline
\textbf{Notation} & \textbf{Convention} \\
\hline
$x$ & Scalar (lowercase italic) \\
$\mathbf{x}$ & Vector (lowercase bold) \\
$X$ & Matrix (uppercase italic) \\
$\mathcal{X}$ & Set or space (calligraphic) \\
$\hat{x}$ & Estimate or empirical value \\
$x^*$ & Optimal or equilibrium value \\
$\bar{x}$ & Reference or baseline value \\
$\tilde{x}$ & Transformed or normalized value \\
$\dot{x}$ & Time derivative: $dx/dt$ \\
$\partial x / \partial y$ & Partial derivative \\
$\nabla$ & Gradient operator \\
$\|\cdot\|$ & Norm (typically Frobenius for matrices) \\
$\mathbb{E}[\cdot]$ & Expectation operator \\
$\mathbb{P}(\cdot)$ & Probability \\
$\mathbf{1}[\cdot]$ & Indicator function \\
$\sim$ & Distributed as \\
$\approx$ & Approximately equal \\
$\propto$ & Proportional to \\
$\in$ & Element of \\
$\subset$ & Subset of \\
$\forall$ & For all \\
$\exists$ & There exists \\
\hline
\end{tabular}
\end{center}

%==============================================================================
\subsection{Subscript and Superscript Conventions}
%==============================================================================

\begin{center}
\renewcommand{\arraystretch}{1.3}
\begin{tabular}{lp{10cm}}
\hline
\textbf{Index} & \textbf{Meaning} \\
\hline
$i, j$ & Component indices (activities, goods, dimensions) \\
$t$ & Time period \\
$d$ & FEPSDE dimension (F, E, P, S, D, Eco) \\
$k$ & Individual or agent index \\
$g$ & Group index \\
$n$ & Nation or country index \\
$^{INU}, ^{KNU}, ^{IDN}$ & Utility category \\
$^{+}, ^{-}$ & Gains, losses (prospect theory) \\
$_{pre}, _{post}$ & Before/after intervention \\
$_{obs}, _{pred}$ & Observed, predicted \\
\hline
\end{tabular}
\end{center}

%==============================================================================
\subsection{Key Equations Reference}
%==============================================================================

For convenience, the most important equations are collected here with chapter references:

\begin{enumerate}
    \item \textbf{Reference Structure Definition} (Ch. 6):
    \begin{equation}
        C^* = \Phi(C^*, \Psi) \quad \text{(fixed point)}
    \end{equation}
    
    \item \textbf{Coherence Index} (Ch. 6):
    \begin{equation}
        K = 1 - \frac{\|C - C^*\|}{\|C^*\|}
    \end{equation}
    
    \item \textbf{Welfare Function} (Ch. 10):
    \begin{equation}
        W = Q \cdot g(K) = Q \cdot g\left(1 - \frac{\|C - C^*\|}{\|C^*\|}\right)
    \end{equation}
    
    \item \textbf{Context-Dependent Reference} (Ch. 9):
    \begin{equation}
        C^*(\Psi) = C^*_0 + \sum_k \beta_k \Psi_k + \sum_{k,l} \gamma_{kl} \Psi_k \Psi_l
    \end{equation}
    
    \item \textbf{Loss Aversion with Context} (Ch. 3):
    \begin{equation}
        \lambda(\Psi) = \lambda_0 + \alpha_C \Psi_C + \alpha_E \Psi_E + \alpha_T \Psi_T
    \end{equation}
    
    \item \textbf{Awareness Function} (Ch. 11):
    \begin{equation}
        A(\Psi) = f(attention, beliefs, processing | \Psi_C, \Psi_K, \Psi_I)
    \end{equation}
    
    \item \textbf{Deviation Diagnostic} (App. AK):
    \begin{equation}
        \Delta = C - C^* \quad \rightarrow \quad \text{Type I-V classification}
    \end{equation}
    
    \item \textbf{Self-Reinforcement Update} (App. AL):
    \begin{equation}
        C^*_{t+1}(\Psi) = C^*_t(\Psi) + \eta \cdot \text{UpdateRule}(\Delta, \text{Type})
    \end{equation}
\end{enumerate}

%==============================================================================

%==============================================================================
% NEW SECTIONS FOR GLOSSARY - EBF UPDATE (January 2026)
%==============================================================================

%------------------------------------------------------------------------------
\subsection{The 5C CORE Framework}
\label{sec:glossary-5c}
%------------------------------------------------------------------------------

The EBF framework is organized around five fundamental questions, each 
answered by a dedicated CORE appendix:

\begin{center}
\renewcommand{\arraystretch}{1.3}
\begin{tabular}{@{}cllll@{}}
\toprule
\textbf{C} & \textbf{Question} & \textbf{Appendix} & \textbf{Output} & \textbf{Symbol} \\
\midrule
WHO & Who has utility? & AAA & Aggregation level & $L \in \{0,\ldots,5\}$ \\
WHAT & What is utility? & C & Utility dimensions & $d \in \text{FEPSDE}$ \\
HOW & How do they interact? & B & Complementarity & $\Gamma, \Lambda$ \\
WHEN & When does context matter? & V & Context vector & $\boldsymbol{\Psi} \in [0,1]^8$ \\
WHERE & Where do numbers come from? & BBB & Parameters & $\Theta, E(\theta)$ \\
\bottomrule
\end{tabular}
\end{center}

%------------------------------------------------------------------------------
\subsection{Parameter Matrices ($\Gamma$, $\Lambda$, $\delta$)}
\label{sec:glossary-matrices}
%------------------------------------------------------------------------------

\begin{center}
\renewcommand{\arraystretch}{1.3}
\begin{longtable}{@{}lp{3cm}p{8cm}@{}}
\toprule
\textbf{Symbol} & \textbf{Name} & \textbf{Definition} \\
\midrule
\endhead

\multicolumn{3}{l}{\textbf{Complementarity Matrices}} \\
\midrule
$\Gamma$ & Utility-Utility Complementarity Matrix & $\Gamma = [\gamma_{dd'}]_{d,d' \in \text{FEPSDE}}$; $6 \times 6$ symmetric matrix (15 unique parameters) capturing how utility dimensions reinforce or offset each other \\
$\gamma_{dd'}$ & Complementarity coefficient & Strength of interaction between dimensions $d$ and $d'$; $\gamma > 0$ means complements, $\gamma < 0$ means substitutes \\
$\bar{\gamma}_{dd'}$ & Baseline complementarity & Global average complementarity at $\Psi_k = 0.5$ for all $k$ \\
\midrule
$\Lambda$ & Context-Context Complementarity Matrix & $\Lambda = [\lambda_{kl}]_{k,l \in \{1,\ldots,8\}}$; $8 \times 8$ symmetric matrix (28 unique parameters) capturing interactions between context dimensions \\
$\lambda_{kl}$ & Context interaction coefficient & Strength of interaction between context dimensions $k$ and $l$ \\
\midrule

\multicolumn{3}{l}{\textbf{Modulation Tensor}} \\
\midrule
$\delta$ & Context Modulation Tensor & $\delta = [\delta_{dd',k}]$; $15 \times 8 = 120$ coefficients capturing how each context dimension $k$ modulates each complementarity $\gamma_{dd'}$ \\
$\delta_{dd',k}$ & Modulation coefficient & Effect of $\Psi_k$ on complementarity $\gamma_{dd'}$; $\delta > 0$ means higher $\Psi_k$ increases complementarity \\
\midrule

\multicolumn{3}{l}{\textbf{Context-Adjusted Values}} \\
\midrule
$\gamma_{dd'}(\boldsymbol{\Psi})$ & Context-adjusted complementarity & $\gamma_{dd'}(\boldsymbol{\Psi}) = \bar{\gamma}_{dd'} + \sum_{k=1}^{8} \delta_{dd',k} \cdot (\Psi_k - 0.5)$ \\
\midrule

\multicolumn{3}{l}{\textbf{Methodological Principles (see Appendix XVIII)}} \\
\midrule
$H_0: \gamma = 0$ & Null hypothesis & Additivity assumption; baseline against which complementarity claims must be tested \\
Exclusion Principle & Methodological stance & Using $\gamma = 0$ as null hypothesis; $\gamma$ identifies where models do NOT apply; maintains falsifiability (Fehr-compatible) \\
Integration Principle & (Problematic) & Using $\gamma$ to explain why elements fit together; risks unfalsifiability; avoided in rigorous EBF applications \\
Disciplinary Power & (Lakatos) & The capacity of a model to exclude observations; source of scientific value \\
Multiple Models Strategy & Methodological refinement & Instead of one null model, work with multiple incompatible models; show EBF rationalizes each individually but not jointly; avoids paradigm conflicts \\
Non-Integrability & Key property & Two models cannot be combined without: identification collapse, logical contradiction, or hidden normative choices; the refined disciplining principle \\
\bottomrule
\end{longtable}
\end{center}

%------------------------------------------------------------------------------
\subsection{Epistemic Status Framework (from Appendix BBB)}
\label{sec:glossary-epistemic}
%------------------------------------------------------------------------------

\begin{center}
\renewcommand{\arraystretch}{1.3}
\begin{longtable}{@{}lp{3cm}p{8cm}@{}}
\toprule
\textbf{Symbol} & \textbf{Name} & \textbf{Definition} \\
\midrule
\endhead

\multicolumn{3}{l}{\textbf{Epistemic Status}} \\
\midrule
$E(\theta)$ & Epistemic status tag & Quality score for parameter $\theta$: $E(\theta) = \frac{\sum_s w_s \cdot q_s \cdot r_s}{\sum_s w_s}$ where $w_s$ = source weight, $q_s$ = quality score, $r_s$ = relevance score \\
$w_s$ & Source weight & Weight assigned to evidence source $s$ based on methodology (RCT: 0.9, Panel: 0.7, LLM-MC: 0.6, etc.) \\
$q_s$ & Quality score & Quality of source $s$ (peer-reviewed replicated: 1.0, working paper: 0.6, etc.) \\
$r_s$ & Relevance score & Relevance of source $s$ to target parameter (same context: 1.0, different context: 0.5, etc.) \\
\midrule

\multicolumn{3}{l}{\textbf{Tier System}} \\
\midrule
Tier 1 & Empirical tier & $E(\theta) \geq 0.9$; direct econometric estimates, suitable for high-stakes policy \\
Tier 2 & LLM-MC tier & $0.7 \leq E(\theta) < 0.9$; synthetic estimates validated against data, suitable for pilots \\
Tier 3 & Axiomatic tier & $E(\theta) < 0.7$; theoretical constraints and placeholders, requires local calibration \\
\midrule

\multicolumn{3}{l}{\textbf{Data Status Tags}} \\
\midrule
\texttt{EMP} & Empirical & Direct estimate from published study \\
\texttt{LLM} & LLM-MC & Synthetic prior from LLM Monte Carlo estimation \\
\texttt{ILL} & Illustrative & Plausible value for demonstration only---DO NOT use in production \\
\texttt{THR} & Theoretical & Derived from axiomatic constraints \\
\midrule

\multicolumn{3}{l}{\textbf{Provenance Tuple}} \\
\midrule
$(\theta, E, \mathcal{S}, t)$ & Provenance tuple & Complete specification: parameter value $\theta$, epistemic status $E$, source set $\mathcal{S}$, timestamp $t$ \\
\bottomrule
\end{longtable}
\end{center}

%------------------------------------------------------------------------------
\subsection{Estimation Methods (from Appendix BBB)}
\label{sec:glossary-methods}
%------------------------------------------------------------------------------

\begin{center}
\renewcommand{\arraystretch}{1.3}
\begin{tabular}{@{}clp{7cm}@{}}
\toprule
\textbf{Method} & \textbf{Name} & \textbf{Description} \\
\midrule
A & Econometric & Direct estimation from panel/experimental data \\
B & LLM-MC & LLM Monte Carlo: synthetic estimation using language models \\
C & Experimental & RCT or quasi-experimental designs \\
D & Meta-Analytic & Bayesian synthesis of multiple sources \\
E & Milgrom Tests & Supermodularity diagnostics for complementarity \\
\bottomrule
\end{tabular}
\end{center}

%------------------------------------------------------------------------------
\subsection{Axiom Index (Complete)}
\label{sec:glossary-axioms}
%------------------------------------------------------------------------------

\subsubsection{Context Axioms (from Appendix V)}

\begin{center}
\small
\renewcommand{\arraystretch}{1.2}
\begin{tabular}{@{}llp{7cm}@{}}
\toprule
\textbf{Axiom} & \textbf{Type} & \textbf{Statement} \\
\midrule
Ψ-EXIST & Foundational & Context affects utility outcomes \\
Ψ-DIM & Foundational & Context is finite-dimensional ($K=8$) \\
Ψ-BOUND & Foundational & Dimensions bounded on $[0,1]$ \\
Ψ-COMPLETE & Structural & 8 dimensions jointly sufficient \\
Ψ-NECESSARY & Structural & Each dimension individually necessary \\
Ψ-ORTH & Structural & Dimensions approximately orthogonal \\
Ψ-MEAS & Structural & Each dimension measurable \\
Ψ-STABLE & Dynamic & Context changes slowly \\
Ψ-ENDO & Dynamic & Context potentially endogenous \\
Ψ-MOD & Modulation & Linear modulation of complementarity \\
Ψ-HETERO & Modulation & Heterogeneous effects across pairs \\
\bottomrule
\end{tabular}
\end{center}

\subsubsection{Estimation Axioms (from Appendix BBB)}

\begin{center}
\small
\renewcommand{\arraystretch}{1.2}
\begin{tabular}{@{}llp{7cm}@{}}
\toprule
\textbf{Axiom} & \textbf{Type} & \textbf{Statement} \\
\midrule
BBB-ADDR & Foundational & Every parameter has an address \\
BBB-EPIST & Foundational & Every parameter has epistemic status \\
BBB-PROV & Foundational & Every parameter has provenance \\
BBB-HIER & Structural & Three-tier hierarchy (Empirical > LLM-MC > Axiomatic) \\
BBB-BAYES & Dynamic & Bayesian updating when new evidence arrives \\
BBB-TRANS & Transparency & All estimation choices documented \\
\bottomrule
\end{tabular}
\end{center}

%------------------------------------------------------------------------------
\subsection{Key Equations Reference (Extended)}
\label{sec:glossary-equations-ext}
%------------------------------------------------------------------------------

\begin{center}
\renewcommand{\arraystretch}{1.5}
\begin{tabular}{@{}lp{6cm}l@{}}
\toprule
\textbf{Name} & \textbf{Equation} & \textbf{Reference} \\
\midrule
Context Vector & $\boldsymbol{\Psi} = (\Psi_I, \Psi_S, \Psi_C, \Psi_K, \Psi_E, \Psi_T, \Psi_M, \Psi_F)$ & Appendix V \\
Normalization & $\Psi_k = \frac{X_k - X_k^{\min}}{X_k^{\max} - X_k^{\min}}$ & Appendix V \\
Context Modulation & $\gamma_{dd'}(\boldsymbol{\Psi}) = \bar{\gamma}_{dd'} + \sum_k \delta_{dd',k}(\Psi_k - 0.5)$ & Appendix V, BBB \\
Epistemic Status & $E(\theta) = \frac{\sum_s w_s q_s r_s}{\sum_s w_s}$ & Appendix BBB \\
Bayesian Update & $\theta' = \frac{n\bar{x} + \kappa\mu_0}{n + \kappa}$ & Appendix BBB \\
Total Parameters & $|\Theta| = 15 + 28 + 120 = 163$ & Appendix BBB \\
\bottomrule
\end{tabular}
\end{center}



\subsection{Cross-Reference to Detailed Treatments}
%==============================================================================

\begin{center}
\renewcommand{\arraystretch}{1.2}
\begin{tabular}{lll}
\hline
\textbf{Topic} & \textbf{Main Text} & \textbf{Appendix} \\
\hline
Complementarity theory & Chapter 5 & Appendix B \\
Reference structure & Chapter 6 & Appendix A \\
Mathematical proofs & Chapter 8 & Appendix D \\
Context dimensions ($\Psi$) & Chapter 9 & Appendix V \\
FEPSDE welfare & Chapter 10 & Appendix C \\
Awareness module & Chapter 11 & Appendix AU \\
Willingness module & Chapter 12 & Appendix AV \\
Empirical validation & Chapter 4, 4.X & Appendix AN \\
Deviation epistemics & --- & Appendix AK \\
Self-reinforcement & --- & Appendix AL \\
Nobel integration & --- & Appendix I \\
Falsifiable predictions & --- & Appendix S \\
Metatheory & --- & Appendix T \\
\hline
\end{tabular}
\end{center}

%==============================================================================
% KEY RESULTS (Template Compliance)
%==============================================================================

% \section{Results} marker for template compliance
\subsection{Key Results: Notation Standards Established}
\label{sec:glossary-results}

\begin{table}[htbp]
\centering
\caption{Summary of Notation Standards in Appendix G}
\label{tab:g-results}
\renewcommand{\arraystretch}{1.2}
\begin{tabular}{@{}llp{6cm}@{}}
\toprule
\textbf{Result} & \textbf{Type} & \textbf{Scope} \\
\midrule
G-R1 & Core Symbols & 50+ mathematical symbols standardized \\
G-R2 & $\Psi$ Dimensions & 8 context dimensions fully specified \\
G-R3 & FEPSDE & 6 welfare dimensions × 144 components \\
G-R4 & Axiom Index & 17+ axioms from V, BBB indexed \\
G-R5 & Cross-Reference & Navigation to all 50+ appendices \\
\bottomrule
\end{tabular}
\end{table}

%==============================================================================
% APPLICATION: WORKED EXAMPLE
%==============================================================================

\subsection{Worked Example: Using the Glossary}
\label{sec:glossary-example}

\textbf{Scenario:} A researcher wants to understand the equation $W = Q \cdot g(K)$ from Chapter 6.

\textbf{Step 1: Symbol Lookup}
\begin{itemize}[nosep]
    \item $W$ = Realized welfare (see Core Mathematical Symbols)
    \item $Q$ = Quality index (welfare potential at equilibrium)
    \item $K$ = Coherence index: $K = 1 - \|C - C^*\|/\|C^*\|$
    \item $g(K)$ = Coherence function mapping $K$ to welfare realization
\end{itemize}

\textbf{Step 2: Cross-Reference}
\begin{itemize}[nosep]
    \item $K$ derivation: See Appendix D (FORMAL-CSTAR)
    \item $C^*$ definition: See Appendix B (CORE-HOW)
    \item Empirical estimation: See Appendix BBB (CORE-WHERE)
\end{itemize}

\textbf{Step 3: Context Dependence}
\begin{itemize}[nosep]
    \item Note that $C^* = C^*(\Psi)$, so coherence is context-dependent
    \item See Appendix V for $\Psi$ dimension details
\end{itemize}

\subsection{Implications for Framework Users}
\label{sec:glossary-implications}

\begin{tcolorbox}[colback=green!5!white, colframe=green!75!black,
    title=\textbf{Practical Implications}]

\textbf{For Researchers:}
\begin{itemize}[nosep]
    \item Always reference this glossary when using EBF notation
    \item Use standardized symbols in publications for consistency
    \item Check axiom index before formulating new propositions
\end{itemize}

\textbf{For Practitioners:}
\begin{itemize}[nosep]
    \item Use cross-reference table to find detailed treatments
    \item FEPSDE dimensions guide welfare measurement
    \item $\Psi$ dimensions inform context assessment
\end{itemize}

\textbf{For Developers:}
\begin{itemize}[nosep]
    \item Symbol table provides variable naming conventions
    \item Estimation methods section guides implementation choices
    \item Epistemic status tags ensure appropriate confidence levels
\end{itemize}

\end{tcolorbox}

%==============================================================================
% BACK MATTER
%==============================================================================

\subsection{Summary}

\begin{tcolorbox}[colback=blue!5!white, colframe=blue!75!black,
    title=\textbf{Appendix G Summary: REF-GLOSSARY}]

\textbf{Core Contribution:}
Central reference ensuring terminological and notational consistency across
the entire EBF framework.

\textbf{Key Components:}
\begin{enumerate}[nosep]
    \item \textbf{Core Symbols:} 50+ mathematical symbols with precise definitions
    \item \textbf{$\Psi$ Dimensions:} Complete specification of 8 context dimensions
    \item \textbf{FEPSDE:} Six-dimensional welfare decomposition (144 components)
    \item \textbf{Axiom Index:} Complete listing of axioms from CORE appendices
    \item \textbf{Cross-References:} Navigation map to all detailed treatments
\end{enumerate}

\textbf{Usage:}
This appendix should be consulted whenever notation questions arise.
All other appendices reference symbols defined here.

\end{tcolorbox}

%==============================================================================
% CRITICAL FOUNDATIONS
%==============================================================================

\subsection{Critical Foundations}
\label{sec:glossary-foundations}

\begin{tcolorbox}[colback=red!5!white, colframe=red!75!black,
    title=\textbf{Critical Foundations: Objections and Responses}]

\textbf{Objection 1: The notation is too complex.}

\textit{``50+ symbols is overwhelming for newcomers.''}

\textbf{Response:} The Quick Reference box provides the essential 10 symbols needed
for basic understanding. The full glossary serves as a reference, not a prerequisite.
Readers should learn symbols incrementally as they encounter them in chapters.

\medskip

\textbf{Objection 2: Why eight $\Psi$ dimensions specifically?}

\textit{``The choice seems arbitrary.''}

\textbf{Response:} The eight dimensions emerge from factor analysis of institutional
variation (Appendix V, Section 2). They are empirically grounded, approximately
orthogonal, and jointly sufficient to capture context effects in behavioral data.

\medskip

\textbf{Objection 3: FEPSDE adds unnecessary complexity.}

\textit{``Why not just use utility?''}

\textbf{Response:} Single-dimensional utility conflates distinct welfare components
with different policy implications. FEPSDE enables targeted interventions: financial
nudges differ from health interventions differ from social programs. The 144-component
structure is complete but most applications use aggregated measures.

\end{tcolorbox}

%==============================================================================
% OPEN ISSUES
%==============================================================================

\subsection{Open Issues and Research Agenda}
\label{sec:glossary-openissues}

\paragraph{Issue 1: Notation Evolution.}
As EBF develops, new concepts may require new symbols. A systematic process
for notation updates while maintaining backward compatibility is needed.

\paragraph{Issue 2: Cross-Framework Translation.}
Other behavioral economics frameworks use different notation. A translation guide
mapping EBF symbols to standard economics notation would aid adoption.

\paragraph{Issue 3: Computational Symbols.}
The glossary focuses on theoretical symbols. A companion computational glossary
covering variable names in code implementations would benefit developers.

\paragraph{Issue 4: Multilingual Terminology.}
German-language terms (e.g., Handlungsbereitschaft) appear alongside English.
Systematic bilingual definitions would aid international collaboration.

%==============================================================================
% REFERENCES
%==============================================================================

\subsection{References for Appendix G}
\label{sec:glossary-references}

\begin{itemize}[nosep]
    \item EBF Master Document: Primary source for all definitions
    \item Appendix B (CORE-HOW): Complementarity matrix $\Gamma$ notation
    \item Appendix C (CORE-WHAT): FEPSDE welfare dimensions
    \item Appendix V (CORE-WHEN): Context $\Psi$ dimension definitions
    \item Appendix BBB (CORE-WHERE): Estimation method notation
    \item Appendix AU (CORE-AWARE): Awareness function notation
    \item Appendix AV (CORE-READY): Willingness function notation
\end{itemize}

\textbf{Full citations:} See Master Bibliography (\texttt{bibliography/bcm\_master.bib})

%==============================================================================
% CROSS-REFERENCE FOOTER
%==============================================================================

\begin{tcolorbox}[colback=gray!10!white, colframe=gray!75!black,
    title=\textbf{Cross-Reference Footer}]

\textbf{This Appendix:} G (REF-GLOSSARY)

\textbf{Depends on:}
\begin{itemize}[nosep]
    \item All CORE appendices: Source of symbol definitions
    \item EBF chapters: Context for notation usage
\end{itemize}

\textbf{Required by:}
\begin{itemize}[nosep]
    \item All EBF appendices: Central notation reference
    \item All EBF chapters: Symbol lookup
    \item External publications: Notation standardization
\end{itemize}

\textbf{Reading Path:}
This appendix is a reference, not sequential reading.
Consult as needed when encountering unfamiliar symbols.

\end{tcolorbox}

%%%%%%%%%%%%%%%%%%%%%%%%%%%%%%%%%%%%%%%%%%%%%%%%%%%%%%%%%%%%%%%%%%%%%%%%%%%%%%%
% END OF APPENDIX G
%%%%%%%%%%%%%%%%%%%%%%%%%%%%%%%%%%%%%%%%%%%%%%%%%%%%%%%%%%%%%%%%%%%%%%%%%%%%%%%
